
 \chapter{Application aux fractions} 
 \begin{itemize}
 \item Simplifier les fractions suivantes : 
 \begin{itemize}
 \it $\frac{1~815}{2~385}$ et $\frac{2~184}{7~560}$.
 \it $\frac{8~613}{9~009}$ et $\frac{12~012}{25~025}$.
 \it $\frac{3~339}{5~880}$ et $\frac{17~226}{18~810}$.
 \it $\frac{6~720}{9~405}$ et $\frac{15~120}{30~576}$.
 \it $\frac{5~292}{8~400}$ et $\frac{9~360}{18~144}$. 
 \it $\frac{5~203}{6~149}$ et $\frac{43~659}{68~607}$.
 \it $\frac{168\times462}{336\times360}$ et $\frac{252\times684}{840\times1~260}$. 
 \it $\frac{507\times 451}{861\times 396}$ et $\frac{315\times 693}{924\times 504}$. 
 \it $\frac{441\times 1~815\times 3~339}{588\times 3~150\times 2~385}$. 
 \it $\frac{756\times 336\times 2~205}{252\times 924\times12~096}$. 
 \end{itemize}
 \item Effectuer les opérations suivantes : 
 \begin{itemize}
 \it $\frac{90}{189}+\frac{45}{84}-\frac{75}{126}$.
 \it $\frac{51}{56}+\frac{88}{231}-\frac{26}{39}$.
 \it $\frac{77}{35}-\frac{80}{66}-\frac{56}{105}$. 
 \it $\frac{91}{52}+\frac9{84}-\frac{55}{105}$.
 \it $\frac{96}{160}+\frac{57}{105}-\frac{39}{77}$. 
 \it $\frac{171}{76} -\frac{161}{147} - \frac{35}{60}$.
 \it $\left(\frac{168}{192}-\frac{50}{112}+\frac{65}{273}\right)\times \frac34$.
 \it $\left(\frac{85}{153}+\frac{21}{60}-\frac{35}{225}\right)\times \frac25$. 
 \it $\left(\frac{18}{72}-\frac{23}{115}+\frac{45}{63}\right)\div \frac27$.
 \it $\left(\frac{100}{275}-\frac{45}{189}+\frac{60}{198}\right)\div\frac3{11}$. 
 \end{itemize}
 \it Trouver une fraction égale à $\frac{52}{117}$ et dont la somme des termes soit égale à $325$.
 \it Trouver une fraction égale à $\frac{44}{121}$ dont la différence des termes soit égale à $357$. 
 \it La somme de deux fractions dont l'une est les $\frac37$ de l'autre est égale à $\frac{40}{21}$. Calculer
 les deux fractions.
 \it Trouver une fraction égale à $\frac{192}{300}$ dont le PGCD des termes soit égal à $13$. 
 \it Trouver une fraction égale à $\frac{132}{385}$ dont
 le PPCM des termes soit $2~100$. 
 \it Trouver deux fractions sachant que leur quotient est égal à $\frac{23}5$ et que leur somme est égale à $\frac{35}3$.
 \it Trouver deux fractions sachant que leur quotient
 est égal à $\frac{17}8$ et que leur différence est égale à $\frac{23}{28}$. 
 \it Trouver les fractions égales à $\frac{126}{192}$ et :\begin{enumerate}
 \item Dont les termes soient inférieur à ceux de la fraction proposée ; 
 \item Dont le numérateur soit inférieur à $400$ et le dénominateur supérieur à $500$ ; 
 \item Dont la différence des termes soit égale à $132$.
 \end{enumerate}
 \it Trouver deux fractions ayant pour numérateur $1$, dont les dénominateurs soient deux nombres entiers consécutifs et comprenant entre elles la fraction $\frac{13}{84}$. 
 \it Deux règles graduées placées côte à côte de façon que leurs origines coïncident mesurent respectivement $1,40$ m et $1,68$ m. La première est partagée en $48$ parties égales et la seconde en $32$ parties égales. Déterminer par leurs numéros les traits de deux graduations qui coïncident. 
 \it Deux pièces d'étoffe mesurent respectivement $\frac{225}8$ de mètre et $\frac{175}{12}$ de mètre. 
 On veut les découper en coupons d'égale longueur. quelle est la plus grande longueur possible pour chacun de ces coupons et combien de coupons chacune de ces pièces fournira-t-elle ? 
 \it Trouver la plus petite fraction dont les quotients par $\frac{28}{45}$ et par $\frac{40}{63}$ soient des nombres entiers. Comparer les termes de la fraction irréductible trouvée au PPCM des numérateurs et au PGCD des dénominateurs des deux fractions proposées. 
 \it Soit les fractions $\frac{63}{22}$ et $\frac{99}{20}$. Trouver la plus grande fraction qui soit contenue un nombre entier de fois dans chacun de ces deux fractions. Comparer la fraction trouvée au PGCD des numérateurs et au PPCM des dénominateurs des deux fractions proposées.
 \it Réduire la fraction $\frac{168}{315}$ à sa plus simple expression. \\ Trouver ensuite toutes les fractions égales à la fraction donnée, et à termes plus petit. Quel est le nombre de ces fractions ?\\
 Déterminer une fraction égale à $\frac{168}{315}$ et dont le numérateur et le dénominateur ont pour somme $8~303$. 
 \it Calculer l'excès de l'unité sur la fraction $\frac8{11}$. \\ On ajoute $5$ à chacun des deux termes de cette fraction (on remarquera que la différence des deux termes ne change pas) : calculer l'excès de l'unité sur la fraction ainsi obtenue. Dire d'après cela si la fraction $\frac8{11}$ augmente ou diminue lorsqu'on ajoute un même 
 nombre à ses deux termes. \\En employant un procédé analogue, dire si la fraction $\frac{15}7$ augmente ou 
 diminue lorsqu'on ajoute un même nombre à ses deux termes.
 \it Un marchand revend à raison de $9$ F le mètre une pièce d'étoffe qu'il a achetée à un prix inconnu. \\
 Il vend une première fois les $\frac38$ de la pièce, une deuxième fois les $\frac35$ du reste et une troisième fois la moitié du nouveau reste. Ces trois ventes ont déjà produit une somme égale au prix d'achat total de la pièce, augmenté de $9$ F. Dans une quatrième vente, le marchand vend le reste de la pièce, et son bénéfice total est $144$ F. \\Calculer la longueur de la pièce et le prix d'achat du mètre. 
 \it Une pièce de tissu de $84$ mètres a été vendue à trois acheteurs. Le premier a eu les $\frac25$ de la pièce ; le second a eu une part égale aux $\frac57$ de la
 part vendue au premier ; le troisième a eu le reste.
 Calculer la longueur du tissu vendu à chaque acheteur.\\
 Le premier a payé $45$ F par mètre, les deux autres ont payé $70$ F par mètre. 
Calculer le prix d'achat de la pièce, sachant que le bénéfice du marchand est le $\frac13$ du prix d'achat.
 \end{itemize}